\textbf{Problem 1.7 (Dujella, 2026).} Let $D_{m,n}(N)$ denote the number of $D(n)$-$m$-tuples with elements in the set $\{1,2,\ldots,N\}$ for a given positive integer $N$.
\begin{enumerate}
\item If $n$ is a perfect square, is there a constant $C(n)$ depending only on $n$, such that $D_{3,n} \sim C(n)N \log N$? Is it true that $C(n) = C(1) = \frac{3}{\pi^2}$?
\item If $n$ is not a perfect square, is there a constant $C(n)$ depending only on $n$, such that $D_{3,n} \sim C(n)N$?
\item If $n$ is a perfect square, what can be said about the asymptotic behaviour of $D_{4,n}$?
\item If $n$ is not a perfect square, is it true that $D_{4,n}$ is finite?
\end{enumerate}

\medskip

\textit{Remark.} An affirmative answer to question (b) in the case when $|n|$ is prime or $n = -1$ is given in Adžaga, N., Dražić, G., Dujella, A., Pethő, A.: Asymptotics of $D(q)$-pairs and triples via $L$-functions of Dirichlet characters. Ramanujan J. \textbf{66}, Article 11 (2025).

Questions (a) and (b) are completely solved in Badesa, J.: On the asymptotics of $D(n)$-pairs and triples. Ramanujan J. \textbf{67}, Article 101 (2025). The answers to both questions are affirmative. In particular, for $n$ a perfect square, $D_{3,n} \sim \frac{3}{\pi^2}N \log N$.

An alternative proof of (a) and (b) is given in the recent paper Dražić, G., Kazalicki, M., Mrazović, R.: Equidistribution of Diophantine pairs among the equivalence classes of quadratic forms. Preprint (2025). arXiv: 2512.22902.
